\documentclass{article}
\usepackage[left=2cm, right=2cm, top=2cm, bottom=2cm]{geometry}
\usepackage{graphicx}
\usepackage{amsmath}
\usepackage{amssymb}
\usepackage{amsthm}
\usepackage{fancyhdr}
\usepackage{verbatim}
\usepackage{listings}
\usepackage{xcolor}
\usepackage{pdfpages}
\usepackage{cancel}
\usepackage{physics}
\usepackage{esint}
\usepackage{braket}

\lstset{
    language=C++,
    basicstyle=\ttfamily\footnotesize,
    keywordstyle=\color{blue},
    commentstyle=\color{green},
    stringstyle=\color{red},
    numbers=left,
    numberstyle=\tiny\color{gray},
    frame=single,
    breaklines=true,
    captionpos=b,
}


\title{MTH 446: Lecture \# 1 (catch up)}
\author{Cliff Sun}

\newtheorem{theorem}{Theorem}[section]
\newtheorem{lemma}[theorem]{Lemma}
\newtheorem{definition}[theorem]{Definition}
\newtheorem{conjecture}[theorem]{Conjecture}
\newtheorem{proposition}[theorem]{Proposition}
\newtheorem{corollary}[theorem]{Corollary}
\newtheorem{one minute paper}[theorem]{One Minute Paper}

\pagestyle{fancy}
\lhead{\textbf{\thepage}\ \ \nouppercase{\rightmark}}
\chead{MTH 446: Lecture \# 1 (catch up)}
\rhead{Cliff Sun}

\begin{document}

\maketitle

\section*{Lecture Span}
\begin{enumerate}
    \item Review of complex numbers
\end{enumerate}

\begin{definition}
    A \underline{set} A is a collection of elements combined. They also follow a rule.
\end{definition}

Some examples:

\begin{enumerate}
    \item The empty set $\empty$ is the set with no elements
    \item If $x$ is an element of $A$, then we write it as $x \in A$. Else, $x \notin A$. 
\end{enumerate}

Following abbreviations, $x \in A \iff$ x is contained in A. Else, x is not contained in A. Some important sets of numbers:

\begin{enumerate}
    \item $\mathbb{N} = \{1,2,3,\dots\}$ (set of Natural numbers)
    \item $\mathbb{Z} = \{0,1,-1,2,-2,\dots\}$ (set of Integers)
    \item $\mathbb{Q} - \{p/q: p,q \in \mathbb{Z}\}$ and $\gcd(p, q) = 1$. (set of rational numbers)
    \item $\mathbb{R}$, the set of reals 
\end{enumerate}

If $A,B \neq \empty$, then $y = f(x)$ for $x \in A, y \in B$ is a mapping which assigns each $x$ to a unique $y$, denoted as $f: A \rightarrow B$, also called a mapping. 

If $A,B \neq \empty$, then \underline{Cartesian product} $A \times B$ is $\{(x,y): x \in A, y \in B\}$

Standard set theory notations for $A,B$ sets:

\begin{enumerate}
    \item $A \subseteq B \iff x \in A \implies x \in B$
    \item $x \in A \cup B \iff (x \in A) \lor (x \in B)$
    \item $x \in A \cap B \iff (x \in A) \land (x \in B)$
    \item $x \in A / B \iff (x \in A) \land (x \notin B)$
\end{enumerate}

Rules of addition:
\begin{enumerate}
    \item Commutative
    \item Associative
    \item Additive identity
    \item Additive inverse
\end{enumerate}

Rules of multiplication
\begin{enumerate}
    \item Commutative
    \item Associative
    \item Multiplicative identity
    \item Multiplicative inverse
    \item Distributive
\end{enumerate}

Let $F$ be a non-empty set on which operations of addition and multiplication are given. Then $F$ is an (algebraic) field if 
\begin{enumerate}
    \item $F$ contains more than one element
    \item $F$ with these properties satisfies all addition and multiplication properties. 
\end{enumerate}

Thus, we say that (F, "+", "$\cdot$") is an Algebraic field. As an example, ($\mathbb{R}$, "+", "$\cdot$") is an Algebraic field in the field of real numbers. 

\section*{Hamilton's construction of complex numbers}

Recall $\mathbb{R^2} = \{(x,y): x,y \in \mathbb{R}\}$ and $z = (x,y)$. In complex analysis, define 

\begin{enumerate}
    \item $\mathbb{C} = \mathbb{R}^2$ (complex plane)
    \item $z = (x,y)$ (complex number)
    \item $x = \Re(z)$ (real part of $z$)
    \item $y = \Im(z)$ (imaginary part of $z$)
\end{enumerate}

Therefore, in $\mathbb{R}^2$ (and therefore in $\mathbb{C}$), we can define addition as 

\begin{equation}
    z_1 + z_2 = (x_1 + x_2, y_1 + y_2)
\end{equation}

Next, introduce multiplication, define 

\begin{equation}
    z_1z_2 = (x_1x_2 - y_1y_2, x_1y_2 + x_2y_1)
\end{equation}

Such operations satisfy all rules of addition and multiplication as well as distribution. 

\section*{Euler's notation}

$z = x + iy$. 

\end{document}